\subsection{等价测度}

命题 10. --- 设 \(\mu\) 和 \(\nu\) 是 T 上的两个正测度。下列条件等价：

\begin{enumerate}
\def\labelenumi{\alph{enumi})}
\item
  局部可忽略集对 \(\mu\) 和对 \(\nu\) 是相同的。
\item
  由 \(\mu\) 和 \(\nu\) 在 \(\mathcal{M}(\mathrm{T})\) 中生成的带是相同的。
\item
  有 \(\nu = g \cdot \mu\) ，其中 \(g\) 是局部 \(\mu\) -可积的，且 \(g(t) > 0\) 关于 \(\mu\) 局部几乎处处成立。
\end{enumerate}

条件 a) 与 b) 依 No.~5 的定理 2 的推论 2 而等价。若它们成立，则 \(\nu = g \cdot \mu\) 且 \(\mu = h \cdot \nu\) ，其中 \(g\) （相应地，\(h\) ）对 \(\mu\) （相应地，\(\nu\) ）是正的且局部可积的。因此（No.~4, 命题 8）\(hg\) 是局部 \(\mu\) -可积的，且 \(\mu = (hg) \cdot \mu\) 。由此推出（No.~3, 命题 3 的推论 2）\(hg\) 关于 \(\mu\) 局部几乎处处等于 1 ，从而 \(g(t) > 0\) 与 \(h(t) = 1/g(t)\) 关于 \(\mu\) 局部几乎处处成立。反之，设 \(\nu = g \cdot \mu\) ，其中 \(g(t) > 0\) 关于 \(\mu\) 局部几乎处处成立；

由于 \((1/g)g\) 局部几乎处处有定义且是局部 \(\mu\) -可积的，1/g 是局部 \(\nu\) -可积的，且 \((1/g)\cdot\nu=\mu\) （No.~4, 命题 8）。

定义 3. --- 局部紧空间 T 上的两个复测度 \(\theta, \theta'\) 称为等价的，如果测度 \(|\theta|\) 和 \(|\theta'|\) 满足命题 10 的条件 a)、b)、c)。

因此，为使 \(\theta\) 和 \(\theta'\) 等价，必须且只须 \(|\theta|\) 和 \(|\theta'|\) 等价。

注. --- 若 \(\mu\) 和 \(\nu\) 是两个等价的正测度，则定义在 T 上、取值于任意拓扑空间 G 中的可测函数对 \(\mu\) 和对 \(\nu\) 是相同的，这由 No.~3 的命题 4 立即得出。

命题 11. --- 设 \(\mu\) 是 T 上的一个正测度。若 T 在无穷远处可数，则存在一个连续函数 \(h\) ，使得 \(h(t) > 0\) 对所有 \(t \in \mathrm{T}\) 成立，且使测度 \(\nu = h \cdot \mu\) （与 \(\mu\) 等价）是有界的。

设 \((\mathrm{K}_{n})\) 是构成 T 的覆盖的一个紧集序列，且对每个 n ，设 \(f_{n}\) 是 \(\mathcal{K}(\mathrm{T})\) 中的一个函数，满足 \(0 \leqslant f_{n} \leqslant 1\) 且 \(f_{n}(t) = 1\) 在 \(\mathrm{K}_{n}\) 上成立（Ch. III, §1, No.~2, 引理 1）。设 \((a_{n})\) 是一个数 \textgreater{} 0 的序列，使得 \(\sum_{n} a_{n} < +\infty\) ；于是级数 \(h = \sum_{n} a_{n} f_{n}\) 在 T 中正规收敛，因而 h 是 T 上的连续函数，使得 \(h(t) > 0\) 对所有 \(t \in \mathrm{T}\) 成立，这是按构造方式得到的。置 \(\nu = h \cdot \mu\) ，我们就有（命题 3 及 Ch. IV, §1, No.~3, 命题 13）

\[
\nu^ {*} (1) = \int^ {*} h d \mu \leqslant \sum_ {n} a _ {n} \int f _ {n} d \mu .
\]

例如，取 \(a_{n} = 2^{-n}\left(\int f_{n}d\mu\right)^{-1}\) 当 \(\int f_n d\mu > 1\) 时，否则取 \(a_{n} = 2^{-n}\) ，我们就有了 \(\sum_{n}a_{n} < +\infty\) 且 \(\nu^{*}(1) < +\infty\) ，这就证明了命题。

命题 12. --- 设 \((\mu_n)\) 是 T 上的一个有界正测度序列；则存在 T 上的一个有界正测度 \(\mu\) ，使得关系 \(\mu^*(\mathrm{N}) = 0\) 等价于 \(\langle \mu_n^*(\mathrm{N}) = 0 \text{ 对每个 } n \rangle\) ；每个测度 \(\mu_n\) 都以 \(\mu\) 为基。此外，若 \(\mu'\) 是另一个具有此性质的 T 上的正测度，则 \(\mu\) 和 \(\mu'\) 等价。

陈述的最后一部分由定义 3 立即得出。为证明 \(\mu\) 的存在性，可以限于 \(\mu_n \neq 0\) （对每个 \(n\) ）的情形；于是测度族 \(\mu_n / 2^n \| \mu_n \|\) 在 \(\mathcal{M}(\mathrm{T})\) 中是可和的，其和 \(\mu\) 满足 \(\| \mu \| \leqslant 1\) 。此外，由于 \(\mu_n \leqslant 2^n \| \mu_n \| \cdot \mu\) ，关系 \(\mu(\mathrm{N}) = 0\) 蕴涵 \(\mu_n(\mathrm{N}) = 0\) 对一切 \(n\) 成立；反之，若 \(\mathrm{N}\) 是一个对所有 \(\mu_n\) 都可忽略的集，则它对 \(\mu\) 是局部可忽略的（§2, No.~2, 命题 1 的推论 2），从而由于 \(\mu\) 有界，它是 \(\mu\) -可忽略的（§1, No.~2, 命题 7 的推论 2）。

